\documentclass[12pt,a4paper]{article}
\usepackage[utf8]{vietnam}
\usepackage{amsmath,amssymb}
\usepackage{geometry}
\geometry{
  top=12mm,
  bottom=12mm,
  left=14mm,
  right=14mm,
  footskip=6mm,
  headheight=14pt
}

\usepackage{tikz}
\usetikzlibrary{arrows.meta,calc,shapes.geometric}
\usepackage[most]{tcolorbox}
\usepackage{fontawesome}
\usepackage{fancyhdr}
\usepackage{enumitem}

\pagestyle{fancy}
\fancyhf{}
\fancyhead[L]{\footnotesize\color{teal!80!black}\faBook\ \textbf{SỔ TAY KIẾN THỨC TOÁN 3}}
\fancyhead[R]{\footnotesize\color{teal!80!black}\textbf{Lớp Toán Thầy Lĩnh -- loptoanthaylinh.vn}}
\fancyfoot[C]{\footnotesize\color{gray}-- Sổ tay kiến thức: Bài 04. Ôn tập bảng nhân 2, bảng nhân 5 --}
\renewcommand{\headrulewidth}{0.4pt}
\renewcommand{\footrulewidth}{0pt}

\setlength{\parindent}{0pt}
\setlength{\parskip}{0pt}

\begin{document}

% --- BANNER BÀI HỌC ---
\begin{tcolorbox}[
  enhanced,
  colback=teal!8!white,
  colframe=teal!70!black,
  arc=3mm, boxrule=1.2pt,
  left=3mm, right=3mm, top=2mm, bottom=2mm,
  halign=center
]
  {\LARGE\bfseries\color{teal!90!black}\faBookmark\ BÀI 04: ÔN TẬP BẢNG NHÂN 2, BẢNG NHÂN 5}\\[1mm]
  {\itshape\footnotesize\color{black!75} Bản chất phép nhân là tổng các số hạng bằng nhau, quy luật bảng nhân 2 và 5 trên 1 trang A4}
\end{tcolorbox}

\vspace{-2mm}

\noindent
\begin{minipage}[t]{0.49\textwidth}\vspace{0pt}
\begin{tcolorbox}[
  enhanced,
  colback=blue!4!white,
  colframe=blue!75!black,
  arc=2mm, boxrule=0.8pt,
  left=2mm, right=2mm, top=1.8mm, bottom=1.8mm,
  title={\footnotesize\bfseries\color{white}\faTimes\ 1. Bảng nhân 2 (Quy luật đếm thêm 2)}
]
{\footnotesize
\begin{itemize}\setlength{\itemsep}{0.8pt}\setlength{\leftskip}{-2.5mm}
  \item \textbf{Bản chất:} $2 \times k$ là lấy $2$ cộng $k$ lần.
  \item \textbf{Quy luật:} Tích của bảng nhân 2 là các số chẵn liên tiếp cách nhau 2 đơn vị ($2, 4, 6, 8, 10, 12, 14, 16, 18, 20$).
  \item Các kết quả đều có chữ số tận cùng là $0, 2, 4, 6, 8$.
  \item \textit{Ví dụ:} $2 \times 6 = 12$; $2 \times 7 = 14$; $2 \times 8 = 16$.
\end{itemize}
}
\end{tcolorbox}
\end{minipage}
\hfill
\begin{minipage}[t]{0.49\textwidth}\vspace{0pt}
\begin{tcolorbox}[
  enhanced,
  colback=teal!4!white,
  colframe=teal!75!black,
  arc=2mm, boxrule=0.8pt,
  left=2mm, right=2mm, top=1.8mm, bottom=1.8mm,
  title={\footnotesize\bfseries\color{white}\faTimes\ 2. Bảng nhân 5 (Quy luật đếm thêm 5)}
]
{\footnotesize
\begin{itemize}\setlength{\itemsep}{0.8pt}\setlength{\leftskip}{-2.5mm}
  \item \textbf{Bản chất:} $5 \times k$ là lấy $5$ cộng $k$ lần.
  \item \textbf{Quy luật:} Tích của bảng nhân 5 tăng đều 5 đơn vị ($5, 10, 15, 20, 25, 30, 35, 40, 45, 50$).
  \item Chữ số tận cùng luôn luôn là \textbf{$0$ hoặc $5$}.
  \begin{itemize}\setlength{\itemsep}{0.5pt}\setlength{\leftskip}{-2.5mm}
    \item Nhân số chẵn: Tận cùng bằng $0$ ($5 \times 4 = 20$).
    \item Nhân số lẻ: Tận cùng bằng $5$ ($5 \times 7 = 35$).
  \end{itemize}
\end{itemize}
}
\end{tcolorbox}
\end{minipage}

\vspace{-2.5mm}
% --- KHỐI 3: SƠ ĐỒ TƯ DUY TIKZ ---
\noindent
\begin{tcolorbox}[
  enhanced,
  colback=orange!4!white,
  colframe=orange!80!black,
  arc=2mm, boxrule=0.8pt,
  left=2mm, right=2mm, top=1.5mm, bottom=1.5mm,
  title={\footnotesize\bfseries\color{white}\faSitemap\ SƠ ĐỒ TƯ DUY: QUY LUẬT DÃY TÍCH CỦA BẢNG 2 VÀ BẢNG 5}
]
\centering
\begin{tikzpicture}[
  box/.style={draw, rounded corners=2pt, align=center, font=\scriptsize\bfseries, inner sep=2.5pt, minimum height=0.75cm},
  arr/.style={-{Stealth[length=2.5mm]}, thick}
]
  \node[box, fill=blue!12, draw=blue!80!black, text width=4.5cm] (n2) at (2.4, 0) {BẢNG NHÂN 2\\Đếm thêm 2 đơn vị\\$2 \to 4 \to 6 \to 8 \to \dots \to 20$};
  \node[box, fill=teal!12, draw=teal!80!black, text width=4.5cm] (n5) at (7.8, 0) {BẢNG NHÂN 5\\Đếm thêm 5 đơn vị\\$5 \to 10 \to 15 \to \dots \to 50$};

  \draw[arr, purple!80!black] (n2) -- (n5) node[midway, above, font=\tiny\bfseries] {Đổi chỗ $2 \times 5 = 5 \times 2$};
\end{tikzpicture}
\end{tcolorbox}

\vspace{-2.5mm}
% --- KHỐI 4: VÍ DỤ & MẸO TRÁNH BẪY ---
\noindent
\begin{minipage}[t]{0.49\textwidth}\vspace{0pt}
\begin{tcolorbox}[
  enhanced,
  colback=green!4!white,
  colframe=green!60!black,
  arc=2mm, boxrule=0.8pt,
  left=2mm, right=2mm, top=1.8mm, bottom=1.8mm,
  title={\footnotesize\bfseries\color{white}\faCheckCircle\ Ví dụ thực hành chuẩn mực}
]
{\footnotesize
\textbf{1. Tính nhẩm:}\\
$2 \times 4 = \mathbf{8}$;\quad $2 \times 9 = \mathbf{18}$.\\
$5 \times 6 = \mathbf{30}$;\quad $5 \times 8 = \mathbf{40}$.\\
\textbf{2. Tính chất giao hoán (đổi chỗ):}\\
$2 \times 5 = 5 \times 2 = \mathbf{10}$.\\
\textbf{3. Bài toán thực tế:}\\
Mỗi bàn có $2$ học sinh. Hỏi $7$ bàn có bao nhiêu học sinh?\\
\textit{Bài giải:} Số học sinh là $2 \times 7 = \mathbf{14}$ (học sinh).
}
\end{tcolorbox}
\end{minipage}
\hfill
\begin{minipage}[t]{0.49\textwidth}\vspace{0pt}
\begin{tcolorbox}[
  enhanced,
  colback=red!4!white,
  colframe=red!70!black,
  arc=2mm, boxrule=0.8pt,
  left=2mm, right=2mm, top=1.8mm, bottom=1.8mm,
  title={\footnotesize\bfseries\color{white}\faExclamationTriangle\ Mẹo tránh bẫy trắc nghiệm}
]
{\footnotesize
\begin{itemize}\setlength{\itemsep}{1.5pt}\setlength{\leftskip}{-2.5mm}
  \item \textbf{Bẫy viết ngược phép nhân:}\\
  $2$ được lấy $7$ lần phải viết là $2 \times 7$, không được viết $7 \times 2$ dù kết quả bằng nhau.
  \item \textbf{Mẹo nhớ nhanh $5 \times k$:}\\
  Nếu $k$ chẵn: Chia đôi $k$ rồi thêm số $0$ vào đuôi (VD: $5 \times 8 \to 8:2=4 \to \mathbf{40}$).
  \item \textbf{Bẫy nhầm phép nhân với phép cộng:}\\
  $2 \times 3 = 6$ (chứ không phải $2 + 3 = 5$).
\end{itemize}
}
\end{tcolorbox}
\end{minipage}


\end{document}
